\chapter*{ВВЕДЕНИЕ}
\addcontentsline{toc}{chapter}{ВВЕДЕНИЕ}

Генерация текста является одной из ключевых задач в области обработки естественного языка, которая активно развивается благодаря достижениям в машинном обучении и увеличению объемов доступных данных. Автоматизированные системы генерации текста находят широкое применение в различных сферах, таких как создание контента, автоматизация общения с пользователями и анализ данных.

Целью данной работы является анализ существующих методов генерации текста, применяемых для создания осмысленного и грамматически правильного текста на русском языке. Для достижения этой цели необходимо решить следующие задачи:
\begin{itemize}
    \item[---] Обозначить основные вехи развития; 
    \item[---] Формализовать задачу генерации текста;
    \item[---] Перечислить методы решения задачи генерации текста;
    \item[---] Сравнить перечисленные методы.
\end{itemize}
